\section{Local structure and two-site blocking}

\begin{definition}[Local simple-MPDO structure]
    \label{def:simple_mpdo_local_structure_data}
    \lean{MPOTensor.SimpleMPDOLocalStructureData}
    \leanok
    \uses{thm:sal_implies_eta_structure}
    This structure records the local ingredients of
    \cite[Appendix~C.2]{Cirac2016MPDO_arXiv}: a normalized three-site reduced
    state with equality in strong subadditivity, the resulting Markov
    decomposition, and a primitive matrix $T$ satisfying
    \[
        \operatorname{tr}(T)=1,
        \qquad
        \operatorname{tr}(T^m)=\operatorname{tr}(T)
    \]
    for every positive integer $m$, together with a rank-one factorization
    \[
        T = a b^\top, \qquad a \cdot b = 1.
    \]
\end{definition}

\begin{theorem}[Local simple-MPDO structure from corrected rank-one criteria]
    \label{thm:simple_mpdo_local_structure_data_of_sal_zcl}
    \lean{MPOTensor.SimpleMPDOLocalStructureData.ofSALZCLOfPairingIdempotent}
    \lean{MPOTensor.SimpleMPDOLocalStructureData.ofSALZCLOfPosSemidef}
    \leanok
    \uses{def:simple_mpdo_local_structure_data, thm:sal_implies_eta_structure,
        thm:psd_trace_powers_rank_one_t, thm:pairing_idempotent_rank_one_t}
    There are two sufficient forms of the local conclusion associated with
    \cite[Lemmas~C.2, C.4, and~C.5]{Cirac2016MPDO_arXiv}. In the first,
    $T$ is positive semidefinite,
    $\operatorname{tr}(T)=1$, and
    $\operatorname{tr}(T^m)=1$ for every positive integer $m$; the spectral
    theorem then gives
    \[
        T = a b^\top, \qquad a \cdot b = 1.
    \]
    In the second, $T$ is the sector trace matrix of linearly independent
    sector tensors satisfying the zero-correlation-length pairing identity.
    This identity implies $T^2=T$, so trace one gives the same rank-one
    factorization. In this second form the constancy of the positive-power
    traces is a conclusion rather than an assumption.
\end{theorem}

\begin{figure}[ht]
    \centering
    \TNMPDOBNTVerticalProduct
    \caption{Orthogonality of the physical supports of a basis of normal
        tensors of a simple renormalization-fixed-point tensor: contracting
        $\mathcal K_y$ with $\mathcal K_x$ through the shared physical index
        vanishes whenever $x\ne y$. This is
        \cite[Theorem~4.9(iv), lines~863--868]{Cirac2016MPDO_arXiv}.}
    \label{fig:mpdo_bnt_vertical_product}
\end{figure}

\begin{figure}[ht]
    \centering
    \TNMPDOProjectorAbsorption
    \caption{Projector absorption for a basis tensor: expanding the identity
        $\sum_jP_j=\Id$ on the ket and bra legs of $\mathcal K_i$ and using
        the local orthogonality relations $P_j\mathcal K_iP_{j'}=0$ for
        $j\ne j'$ or $i\ne j$ leaves only the diagonal term, so
        $\mathcal K_i=P_i\mathcal K_i=\mathcal K_iP_i$. This is the display
        in the proof of
        \cite[Appendix~C.2, lines~1745--1751]{Cirac2016MPDO_arXiv}, and it
        implies the orthogonality of
        Figure~\ref{fig:mpdo_bnt_vertical_product}.}
    \label{fig:mpdo_projector_absorption}
\end{figure}

\begin{figure}[ht]
    \centering
    \TNMPDOTwoSiteBlocking
    \caption{The two-site blocking of the simple tensor $K$: the blocked
        tensor $M$ is obtained by contracting two copies of $K$ along one
        virtual bond and grouping the two ket legs and the two bra legs into
        blocked physical indices. This is
        \cite[Theorem~4.9, lines~851--856]{Cirac2016MPDO_arXiv}.}
    \label{fig:mpdo_two_site_blocking}
\end{figure}

\begin{definition}[Two-site blocking of an MPO tensor]
    \label{def:mpdo_two_site_blocking}
    \lean{MPOTensor.blockTwo}
    \leanok
    For an MPO tensor $K$, its two-site blocking $K^{[2]}$ has physical
    indices $(i_0,i_1)$ and $(j_0,j_1)$ and matrices
    \[
        \bigl(K^{[2]}\bigr)^{(i_0,i_1),(j_0,j_1)}
          =K^{i_0j_0}K^{i_1j_1}.
    \]
    This is the blocking used in
    \cite[Theorem~4.9, lines~851--856]{Cirac2016MPDO_arXiv}.
\end{definition}

\begin{definition}[Four-site physical closure]
    \label{def:mpdo_four_site_physical_closure}
    \lean{MPOTensor.finFourArrowEquiv, MPOTensor.physClose4}
    \leanok
    For a virtual matrix $X$, the four-site physical closure of $K$ is the
    operator $K_4(X)$ whose matrix coefficients are
    \[
        K_4(X)_{(i_0,(i_1,(i_2,i_3))),(j_0,(j_1,(j_2,j_3)))}
          =\tr\!\left(
            K^{i_0j_0}K^{i_1j_1}K^{i_2j_2}K^{i_3j_3}X\right).
    \]
\end{definition}

\begin{theorem}[General and right-associated four-site closures]
    \label{thm:mpdo_phys_close_four}
    \lean{MPOTensor.physCloseN_four_eq_physClose4}
    \leanok
    \uses{def:phys_close, def:mpdo_four_site_physical_closure}
    Under the canonical right-associated identification of four-site
    configurations with quadruples of physical indices,
    \[
        K_4(X)_{(i_0,(i_1,(i_2,i_3))),(j_0,(j_1,(j_2,j_3)))}
          =\tr\!\left(K^{i_0j_0}K^{i_1j_1}K^{i_2j_2}K^{i_3j_3}X\right).
    \]
\end{theorem}

\begin{proof}\leanok
    Expand the general four-site closure in the right-associated coordinates.
\end{proof}

\begin{theorem}[One-site closure of the two-site blocking]
    \label{thm:mpdo_one_site_closure_two_site_blocking}
    \lean{MPOTensor.physClose1_blockTwo_eq_physClose2}
    \leanok
    \uses{def:phys_close, def:mpdo_two_site_blocking}
    After identifying one blocked physical index with two original indices,
    \[
        \bigl(K^{[2]}\bigr)_1(X)=K_2(X).
    \]
\end{theorem}

\begin{proof}\leanok
    Expand one blocked tensor matrix and decode its two physical indices.
\end{proof}

\begin{theorem}[Two-site closure of the two-site blocking]
    \label{thm:mpdo_two_site_closure_two_site_blocking}
    \lean{MPOTensor.physClose2_blockTwo_eq_physClose4}
    \leanok
    \uses{def:phys_close, def:mpdo_two_site_blocking,
      def:mpdo_four_site_physical_closure}
    After identifying two blocked physical indices with four original
    indices, one has
    \[
        \bigl(K^{[2]}\bigr)_2(X)=K_4(X).
    \]
\end{theorem}

\begin{proof}\leanok
    Expand the two blocked tensor matrices and reassociate their product.
\end{proof}
